\documentclass{beamer}
\usetheme{metropolis}

\usepackage{mathtools}
\usepackage[T1]{fontenc}
\usepackage[utf8]{inputenc}
\usepackage[backend=biber, style=numeric, sorting=ynt]{biblatex}
\addbibresource{refs.bib}
\usepackage{tabularx}
\usepackage{hyperref}
\usepackage{amssymb}
\usepackage{amsmath}
\usepackage{minted}
\usepackage{svg}
\usepackage{twemojis}

\title{Graph drawing by stress minimization}
\author{Laslo Hauschild, Lennart Peters}
\date{\today}

\begin{document}

\begin{frame}
    \titlepage
\end{frame}

\section{Introduction}

\begin{frame}{Introduction}
    Graphs are a prominent object in modern mathematics, useful for abstracting complex reallife situations.
    \begin{itemize}
        \item Railway networks, data structures, geometric objects
        \item Model relationships: vertices for entities, edges for connections
    \end{itemize}
\end{frame}

\begin{frame}
    \textbf{Example:} Two companies with employees as vertices and daily interactions as edges.
    \begin{itemize}
        \item Compare communication hierarchy and department interactions
        \item Difficulty: meticulously searching through huge lists of vertices and edges
    \end{itemize}
    \textbf{Goal:} Represent finite graphs in $\mathbb{R}^2$ and $\mathbb{R}^3$ esthetically.
    \begin{itemize}
        \item Drawing: vertices as points, edges as straight-line segments
        \item Quickly understand structure and type of a given graph
    \end{itemize}
\end{frame}

\section{Graph Embedding}

\begin{frame}{Graph embedding}
    \textbf{Graph definition:}
    \begin{itemize}
        \item set of vertices $V$, set of edges $E \subseteq \{ \{ u,v \} \mid u,v \in E\}$ then $G=(V,E)$
    \end{itemize}
    Our goal is to draw graphs, for our purposes a drawable graph has a finite set of vertices.
    \begin{itemize}
        \item finite $V \implies$ finite $E$
    \end{itemize}
    A graph embedding $p$ is an injective map $p : V \to A$, which maps every $v \in V$. $p$ induces an embedding of $E$, therefore $p$ is a lossless embedding of $G$ in $A$.
    \begin{itemize}
        \item $\exists p \iff |V| \leq |A|$
    \end{itemize}
    
    Let us call $H=p(G)$ and choose $A= \mathbb{R}^n, \quad n=2,3$
\end{frame}

\begin{frame}{Modeling stress}
    We add a drawing function $d : F \to \mathbb{R}, \{ a,b \} \mapsto \overline{ab}$, giving us a drawing of $H$.
    
    $W \subseteq \mathbb{R}^n$ essentially randomly chosen, drawings of $H$ are mostly tangled.
    We use a stress function $S$ that assigns the graph $H$ a real number, the stress in $H$.
    Minima of $S$ are our desired representations of $G$, if possible the global minimum as ideal.
\end{frame}

\begin{frame}{FR model}
    developed 1991, by Fruchterman and Reingold, following two main principles~\cite{fr_paper}:
    \begin{itemize}
    \item drawing neighboring vertices near each other
    \item vertices should not  be drawn too near to each other
    \end{itemize}

    The FR model uses the following stress-function:
    \begin{equation}
        S_{FR}(p(G)) = \frac{1}{6}\sum_{\{u,v\} \in E}{\parallel p(u)-p(v) \parallel^3}
        - \sum_{\{u,v \} \subseteq V}{\log \parallel p(u)-p(v) \parallel}
    \end{equation}
\end{frame}

\begin{frame}
    \textbf{Observe:}
    \begin{itemize}
        \item the further apart neighboring vertices get mapped the bigger the stress
        \item small distances between vertices let the logarithm explode
        \item the logarithm dominates at small distances and hinders the graph from collapsing, or converging into a single point
    \end{itemize}
    $S_{FR}$ satisfies the two main principles, while being rather quick when implemented into an algorithm.
\end{frame}

\begin{frame}{ERLL model}
    The Edge Repulsion LinLog model by Andreas Noack, 2003, places its focus on the graph's edges.
    For $v \in V$, let $\deg(v)$ denote the number of edges $v$ is an endpoint of~\cite{erll_paper}.
    \begin{equation}
    \begin{split}
        S_{ERLL}(p(G)) = & \frac{1}{2} \sum_{\{u,v\} \in E}{\parallel p(u)-p(v) \parallel} \\
        & - \sum_{\{u,v \} \subseteq V}{\deg(u)\deg(v) \log\parallel p(u)-p(v) \parallel}
    \end{split}
    \end{equation}
\end{frame}

\begin{frame}
    \textbf{Observe:}
    \begin{itemize}
        \item edges cause both attraction and repulsion
        \item $G$ of uniform degrees: FR $\approx$ ERLL scaled with $c=\deg{(v)}^2$
        \item $G$ of non-uniform degrees: FR repulsion remains the same, but vertices $u,v$ repel each other proportional to $\deg(u)\deg(v)$ in ERLL.
    \end{itemize}
    $\implies$ ERLL spreads out "big nodes" compared to FR, usefull in highly non-uniform graphs
\end{frame}
\section{Implementation}

\begin{frame}[fragile]{Plotting a Graph Embedding}
\begin{minted}[fontsize=\footnotesize]{julia}
using Plots
using Graphs

function plot_graph(g::SimpleGraph, p::Matrix{Float64})
    plt = plot(legend = false, aspect_ratio = :equal)
    for e in edges(g)
        plot!(plt, [p[1, src(e)], p[1, dst(e)]],
            [p[2, src(e)], p[2, dst(e)]], color = :black)
    end
    scatter!(plt, p[1, :], p[2, :], markersize=5, color = :blue)
    return plt
end
\end{minted}
\end{frame}

\begin{frame}
\begin{itemize}
    \item Dependencies: Graphs~\cite{graphsjl} and Plots~\cite{plotsjl} libraries
    \item Arguments: A graph $g$ and a position matrix $p$ with vertex coordinates
    \item Order of plotting matters: edges first, then vertices
\end{itemize}
\end{frame}

\begin{frame}[fragile]{Plotting a Graph Embedding in 3 Dimensions}
\begin{minted}[fontsize=\footnotesize]{julia}
function plot_graph_3d(g::SimpleGraph, p::Matrix{Float64})
    plt = plot(legend = false, aspect_ratio = :equal)
    for e in edges(g)
        plot!(plt, [p[1, src(e)], p[1, dst(e)]],
            [p[2, src(e)], p[2, dst(e)]],
            [p[3, src(e)], p[3, dst(e)]], color = :black)
    end
    scatter!(plt, p[1, :], p[2, :], p[3, :],
        markersize=5, color = :blue)
    return plt
end
\end{minted}
\end{frame}

\begin{frame}{Stress functions}
    \begin{itemize}
        \item ERLL Stress and FR Stress are structured similarly
        \begin{itemize}
            \item edge stress
            \item vertex stress
            \item combine function
        \end{itemize}
        \item Idea: One function that takes anonymous functions
    \end{itemize}
\end{frame}


\begin{frame}[fragile]{Stress Function}
\begin{minted}[fontsize=\footnotesize]{julia}
function stress(g::SimpleGraph, edge_stress_f::Function, 
    vertex_stress_f::Function, combine_f::Function)

    n = nv(g)
    edge_stress = 0
    for e in edges(g)
        edge_stress += edge_stress_f(src(e), dst(e))
    end
    vertex_stress = 0
    for i in 1:n
        for j in (i+1):n
            vertex_stress += vertex_stress_f(i, j)
        end
    end
    return combine_f(edge_stress, vertex_stress)
end
\end{minted}
\end{frame}

\begin{frame}[fragile]{Stress Function}
\begin{minted}[fontsize=\scriptsize]{julia}
function fr_stress(g::SimpleGraph, p::Matrix{Float64})
    D = pairwise(Euclidean(), p, dims=2)
    edge_stress_f = (i, j) -> D[i, j]^3
    vertex_stress_f = (i, j) -> log(D[i, j] + 10^(-6))
    combine_f = (x, y) -> x / 6 - y
    return stress(g, edge_stress_f, vertex_stress_f, combine_f)
end

function erll_stress(g::SimpleGraph, p::Matrix{Float64}, deg::Vector{Int})
    D = pairwise(Euclidean(), p, dims=2)
    edge_stress_f = (i, j) -> D[i, j]
    vertex_stress_f = (i, j) -> deg[i] * deg[j] * log(D[i, j] + 10^(-6))
    combine_f = (x, y) -> x / 2 - y
    return stress(g, edge_stress_f, vertex_stress_f, combine_f)
end
\end{minted}
\end{frame}

\begin{frame}
\begin{itemize}
    \item We do $\log(D[i, j] + 10^{-6})$ to limit the size of the logarithm
    \item For efficiency, \texttt{pairwise} function (Distances library~\cite{distancesjl}) to compute the distance matrix $D$.
    \item ERLL stress has additional argument: degree of each vertex (why: could calculate in stress function, this would be calculated every time, which is inefficient)
\end{itemize}
\end{frame}

\section{Gradient Descent}

\begin{frame}{Zygote}
    \begin{itemize}
        \item We do not derive the gradients manually
        \item Zygote library~\cite{zygotejl} to compute the gradients automatically
        \item Uses source-to-source automatic differentiation:
        \begin{enumerate}
            \item Transforms code into differentiable form
            \item Applies the chain rule
            \item Done \twemoji{party popper}
        \end{enumerate}
    \end{itemize}
\end{frame}

\begin{frame}{Gradient Descent}
    Principle of gradient descent:
    \begin{enumerate}
        \item Initialize vertex positions randomly in 2D space
        \item Compute stress and its gradient at each step
        \item Update positions minimally in the direction of steepest descent
        \item Check for convergence: stop after max steps or if stress values do not change meaningfully
        \begin{enumerate}[$\rightarrow$]
            \item Do not check for equality, check if relative difference is smaller than a small threshold
        \end{enumerate}
    \end{enumerate}
\end{frame}

\begin{frame}[fragile]
\begin{minted}[fontsize=\scriptsize]{julia}
function simple_gradient_descent(g::SimpleGraph, stress_fn::Function, 
    alpha::Float64, max_steps::Int64, ...)

    n = nv(g) # Number of vertices
    p = rand(2, n) # Initial positions
    stresses = Float64[]

    for k in 1:max_steps
        s, (grad,) = Zygote.withgradient(q -> stress_fn(g, q), p)
        push!(stresses, s)

        if length(stresses) > 1 &&
            abs((stresses[end] - stresses[end-1]) / stresses[end-1]) < 1e-7
            ... # result plotting code
            break
        end

        p -= alpha * grad

        ... # continuous plotting code
    end
end
\end{minted}
\end{frame}

\begin{frame}{Break condition}
    \textbf{Definitons}
    \begin{itemize}
        \item $s_l :=$ last stress value
        \item $s_{l-1} :=$ second to last stress value
    \end{itemize}

    \textbf{Convergence criterion}

    Natural way (absolute difference): $\parallel s_l - s_{l-1} \parallel < \epsilon$.

    \begin{enumerate}[$\rightarrow$]
        \item Problem: Depends on absolute size of stress values.
    \end{enumerate}

    Better way (relative difference): $\parallel \frac{s_l - s_{l-1}}{s_{l-1}} \parallel < \epsilon$.
\end{frame}

\begin{frame}{Nesterov's accelerated gradient descent}
    \textbf{Motivation:} Simple gradient descent can be slow, especially in areas where the stress function is flat.

    \textbf{Idea:} A ball rolling downhill
    \begin{itemize}
        \item Has inertia
        \item Gains speed as it goes down the hill
        \item Can escape local minima and find better solutions
    \end{itemize}

    \textbf{Implementation:} 
    \begin{itemize}
        \item In addition to current gradient, we also consider previous update (momentum) when updating positions
        \item Friction modelled by $\beta = 0.9$, meaning we keep 90\% of previous momentum
    \end{itemize}
\end{frame}

\begin{frame}[fragile]
\begin{minted}[fontsize=\scriptsize]{julia}
function nesterov_gradient_descent(g::SimpleGraph, stress_fn::Function, alpha::Float64,
    beta::Float64, max_steps::Int64, ...)

    n = nv(g) # Number of vertices
    p = rand(2, n) # Initial positions
    stresses = Float64[]
    momentum = zeros(2, n)
    for k in 0:max_steps
        s, (grad,) = Zygote.withgradient(q -> stress_fn(g, q), p + beta * momentum)
        new_momentum = beta * momentum - alpha * grad
        push!(stresses, s)
        if length(stresses) > 50 &&
            abs((stresses[end] - stresses[end-1]) / stresses[end-1]) < 1e-7
            ... # result plotting code
            break
        end
        p += new_momentum
        momentum = new_momentum
        ... # continuous plotting code
    end
end
\end{minted}
\end{frame}

\section{Results}

\begin{frame}{Good convergence}
\begin{figure}[H]
\centering
\includesvg[width=0.7\textwidth]{../assets/wheel_graph_11_simple_fr_a0.001_i4597.svg}
\caption{\label{wheel-graph-good} Wheel Graph 11, Simple Gradient Descent, $\alpha=0.001$, Took 4597 Iterations}
\end{figure}
\end{frame}

\begin{frame}{Bad convergence}
    Less than $10\%$ of the runs, the graphs converge poorly:
\begin{figure}[H]
\centering
\includesvg[width=0.7\textwidth]{../assets/wheel_graph_11_simple_fr_a0.001_i3669.svg}
\caption{\label{wheel-graph-bad} Wheel Graph 11, Simple Gradient Descent, $\alpha=0.001$, Took 3669 Iterations}
\end{figure}
\end{frame}

\begin{frame}{Figuring out good learning rates}
    We had problems with the suggested learning rates.
    \begin{itemize}
        \item FR stress: learning rates were fine
        \item ERLL: learning rates were too small, making convergence very slow
        \item We try different learning rates and see how they affect convergence speed and quality of the results
    \end{itemize}
\end{frame}

\begin{frame}{Learning rate $\alpha=0.01$}
\begin{figure}[H]
\centering 
\includesvg[width=0.7\textwidth]{../assets/wheel_graph_11_simple_fr_a0.01_i637.svg}
\caption{\label{wheel-graph-alpha-0.01} Wheel Graph 11, Simple Gradient Descent, $\alpha=0.01$, Took 428 Iterations}
\end{figure}
\end{frame}

\begin{frame}{Learning rate $\alpha=0.1$}
\begin{figure}[H]
\centering
\includesvg[width=0.7\textwidth]{../assets/wheel_graph_11_simple_fr_a0.1_i58.svg}
\caption{\label{wheel-graph-alpha-0.1} Wheel Graph 11, Simple Gradient Descent, $\alpha=0.1$, Took 57 Iterations}
\end{figure}
\end{frame}

\begin{frame}{Learning rate $\alpha=1.0$}
\begin{figure}[H]
\centering
\includesvg[width=0.7\textwidth]{../assets/wheel_graph_11_simple_fr_a1.0_i11.svg}
\caption{\label{wheel-graph-alpha-1.0} Wheel Graph 11, Simple Gradient Descent, $\alpha=1.0$, Took 11 Iterations}
\end{figure}
\end{frame}

\begin{frame}
Increasing the learning rate\dots
\begin{itemize}
    \item makes convergence faster
    \item increases the rate of poorly converged graphs (difficult to quantify)
    \item makes the stress go up at bigger values
    \item will prevent convergence at very big values because of number overflow
\end{itemize}
\end{frame}


\begin{frame}
    With simple gradient descent and ERLL and $\alpha=0.001$, convergence is painfully slow. Increasing it to $\alpha=0.1$ works:
    
    \begin{figure}[H]
    \centering
    \includesvg[width=0.7\textwidth]{../assets/wheel_graph_11_simple_erll_a0.1_i20994.svg}
    \caption{\label{wheel-graph-erll} Wheel Graph 11, Simple Gradient Descent, $\alpha=0.1$, Took 20994 Iterations}
    \end{figure}
\end{frame}


\begin{frame}{A good learning rate}
    \textbf{FR Stress}
    $\alpha=0.001$ offers reasonable convergence speed and good results.

\textbf{ERLL stress}
$\alpha=0.001$ is too small, making convergence very slow.
Increasing it to $\alpha=0.1$ works.

We will use these learning rates for the rest of the project.
\end{frame}

\begin{frame}{Icosahedral -- ERLL}
    \begin{figure}[H]
    \centering
    Nesterov Gradient Descent, ERLL, $\alpha=0.1$, $\epsilon=1e-7$
    \includesvg[width=0.6\textwidth]{../assets/icosahedral_nesterov_erll_a0.1_i2593.svg}
    \caption{\label{icosahedral-nesterov-erll} Icosahedral Graph, Took 2593 Iterations}
    \end{figure}
\end{frame}

\begin{frame}{Icosahedral -- FR}
    \begin{figure}[H]
    \centering
    Nesterov Gradient Descent, FR, $\alpha=0.001$, $\epsilon=1e-7$
    \includesvg[width=0.6\textwidth]{../assets/icosahedral_nesterov_fr_a0.001_i389.svg}
    \caption{\label{icosahedral-nesterov-fr} Icosahedral Graph, Took 389 Iterations}
    \end{figure}
\end{frame}

\begin{frame}{Karate Network -- ERLL}
    \begin{figure}[H]
    \centering
    Nesterov Gradient Descent, ERLL, $\alpha=0.1$, $\epsilon=1e-7$
    \includesvg[width=0.6\textwidth]{../assets/karate_nesterov_erll_a0.1_i9828.svg}
    \caption{\label{karate-nesterov-erll} Karate Network, Took 9828 Iterations}
    \end{figure}
\end{frame}

\begin{frame}{Karate Network -- FR}
    \begin{figure}[H]
    \centering
    Nesterov Gradient Descent, FR, $\alpha=0.001$, $\epsilon=1e-7$
    \includesvg[width=0.6\textwidth]{../assets/karate_nesterov_fr_a0.001_i2238.svg}
    \caption{\label{karate-nesterov-fr} Karate Network, Took 2238 Iterations}
    \end{figure}
\end{frame}

\begin{frame}{Dorogovtsev Mendes Graph -- ERLL}
    \begin{figure}[H]
    \centering
    Nesterov Gradient Descent, ERLL, $\alpha=0.1$, $\epsilon=1e-7$
    \includesvg[width=0.6\textwidth]{../assets/dorogovtsev_mendes_nesterov_erll_a0.1_i4907.svg}
    \caption{\label{dorogovtsev-mendes-nesterov-erll} Dorogovtsev Mendes Graph, Took 4907 Iterations}
    \end{figure}
\end{frame}

\begin{frame}{Dorogovtsev Mendes Graph -- FR}
    \begin{figure}[H]
    \centering
    Nesterov Gradient Descent, FR, $\alpha=0.001$, $\epsilon=1e-7$
    \includesvg[width=0.6\textwidth]{../assets/dorogovtsev_mendes_nesterov_fr_a0.001_i3335.svg}
    \caption{\label{dorogovtsev-mendes-nesterov-fr} Dorogovtsev Mendes Graph, Took 3335 Iterations}
    \end{figure}
\end{frame}

\begin{frame}{Circular Ladder Graph -- ERLL}
    \begin{figure}[H]
    \centering
    Nesterov Gradient Descent, ERLL, $\alpha=0.1$, $\epsilon=1e-7$
    \includesvg[width=0.6\textwidth]{../assets/circular_ladder_nesterov_erll_a0.1_i4944.svg}
    \caption{\label{circular-ladder-nesterov-erll} Circular Ladder Graph, Took 4944 Iterations}
    \end{figure}
\end{frame}

\begin{frame}{Circular Ladder Graph -- FR}
    \begin{figure}[H]
    \centering
    Nesterov Gradient Descent, FR, $\alpha=0.001$, $\epsilon=1e-7$
    \includesvg[width=0.6\textwidth]{../assets/circular_ladder_nesterov_fr_a0.001_i1292.svg}
    \caption{\label{circular-ladder-nesterov-fr} Circular Ladder Graph, Took 1292 Iterations}
    \end{figure}
\end{frame}

\begin{frame}{Watts Strogatz Graph -- ERLL}
    \begin{figure}[H]
    \centering
    Nesterov Gradient Descent, ERLL, $\alpha=0.1$, $\epsilon=1e-7$
    \includesvg[width=0.6\textwidth]{../assets/watts_strogatz_nesterov_erll_a0.1_i7710.svg}
    \caption{\label{watts-strogatz-nesterov-erll} Watts Strogatz Graph, Took 7710 Iterations}
    \end{figure}
\end{frame}

\begin{frame}{Watts Strogatz Graph -- FR}
    \begin{figure}[H]
    \centering
    Nesterov Gradient Descent, FR, $\alpha=0.001$, $\epsilon=1e-7$
    \includesvg[width=0.6\textwidth]{../assets/watts_strogatz_nesterov_fr_a0.001_i791.svg}
    \caption{\label{watts-strogatz-nesterov-fr} Watts Strogatz Graph, Took 791 Iterations}
    \end{figure}
\end{frame}

\begin{frame}{Static Scale Free Graph -- ERLL}
    \begin{figure}[H]
    \centering
    Nesterov Gradient Descent, ERLL, $\alpha=0.1$, $\epsilon=1e-7$
    \includesvg[width=0.6\textwidth]{../assets/static_scale_free_nesterov_erll_a0.1_i2887.svg}
    \caption{\label{static-scale-free-nesterov-erll} Static Scale Free Graph, Took 2887 Iterations}
    \end{figure}
\end{frame}

\begin{frame}{Static Scale Free Graph -- FR}
    \begin{figure}[H]
    \centering
    Nesterov Gradient Descent, FR, $\alpha=0.001$, $\epsilon=1e-7$
    \includesvg[width=0.6\textwidth]{../assets/static_scale_free_nesterov_fr_a0.001_i884.svg}
    \caption{\label{static-scale-free-nesterov-fr} Static Scale Free Graph, Took 884 Iterations}
    \end{figure}
\end{frame}

\section{Conclusion}

\begin{frame}
\begin{itemize}
    \item FR stress is much faster than ERLL stress
    \item FR stress produces better results than ERLL stress
    \item ERLL stress pushes vertices too far apart
    \item Some graphs get stuck in suboptimal positions, e.g.\ circular ladder graph
    \item Bipartite graphs get drawn in a very unusual way with this method
\end{itemize}

For a truly good graph drawing algorithm, many improvements are possible.
\end{frame}

\begin{frame}[allowframebreaks]{References}
    \printbibliography[heading=none] % 'none' because the Frame Title acts as the heading
\end{frame}

\end{document}